% ============================================================================
% WVC 2026 (Workshop de Visão Computacional) — UNICAMP
% Per the official CFP (ime.unicamp.br/~wvc/main-track.html): 6-page limit,
% IEEE conference template, double-blind review, English title+abstract
% (abstract <=150 words), PDF via CMT3.
%
% DOUBLE-BLIND NOTICE: author names/affiliations below are placeholders.
% Do not restore real identities (or the DK/BR collaboration note in the
% Introduction) until the camera-ready version after acceptance, since the
% CFP explicitly requires no identifying information in the blind
% submission, including in acknowledgments or linked material.
%
% Preamble reconciled 2026-07 against the official IEEE conference template
% found by the author (paper/template-wvc/IEEE-conference-template-062824.tex,
% version 062824), the authoritative source; earlier verification used an
% Overleaf mirror since direct ieee.org access blocked automated fetches.
% Requires IEEEtran.cls, bundled with most TeX distributions and with
% Overleaf's IEEE templates; not vendored in this repo. a4paper matches
% the Brazilian venue; swap to US letter only if the CFP later says otherwise
% (the official template itself is silent on paper size). Kept booktabs
% (for \toprule/\midrule/\bottomrule table rules) and algorithm+algpseudocode
% (for the \State/\If/\ElsIf syntax in Algorithm 1) instead of the official
% template's plain {algorithmic}, since both are standard, widely accepted
% substitutions with IEEEtran and rewriting Algorithm 1 for the older
% package syntax is unnecessary risk; added the official template's
% {cite}, {amssymb}, {amsfonts} for completeness.
%
% NOTE: this paper is deliberately structured around an analysis question
% (does the necessity of audio-visual fusion vary across deployment
% conditions?) rather than around the system itself, and does not cite or
% reference the companion CONICT extended-abstract report on the same
% system per the author's explicit instruction. Keep the System Overview
% section compact; the stratified analysis in Results is this paper's
% primary contribution.
% ============================================================================
\documentclass[conference,a4paper]{IEEEtran}
\IEEEoverridecommandlockouts

\usepackage[utf8]{inputenc}
\usepackage[T1]{fontenc}
\usepackage{cite}
\usepackage{amsmath}
\usepackage{amssymb}
\usepackage{amsfonts}
\usepackage{graphicx}
\usepackage{textcomp}
\usepackage{xcolor}
\usepackage{booktabs}
\usepackage{algorithm}
\usepackage{algpseudocode}
\usepackage{hyperref}

\title{When Is Voice Not Enough? A Deployment-Scale Analysis of\\
Multimodal Necessity in Real-Time Speaker Attribution}

% TODO (camera-ready only): restore real author list / affiliation / emails.
% Blind-submission placeholder below.
\author{
\IEEEauthorblockN{Anonymous Author(s)}
\IEEEauthorblockA{Affiliation withheld for double-blind review}
}

\begin{document}
\maketitle

\begin{abstract}
Multimodal speaker attribution systems are typically evaluated with a
single aggregate accuracy or coverage figure, which implicitly assumes
that the benefit of fusing voice, face, and active-speaker cues is
uniform across deployment conditions. We test this assumption using
deployment logs from a real-time meeting-tracking system that
attributes each speech segment through an explicit, priority-ordered
audio-visual fusion cascade rather than a single learned model. Across
93 real, unscripted meeting sessions and 1,669 speech segments
collected over 2.5 months, we find that reliance on voice biometrics
alone is far from uniform, ranging from 5\% to 40\% depending on
expected group size, and that an apparent decline in voice-only
resolution over the deployment is largely explained by a shift in
session composition rather than by degrading voice quality. The
cascade sustains real-time throughput (89.2\% of segments faster than
real time) and flags its own uncertainty on 28.7\% of segments without
labeled data.
\end{abstract}

\begin{IEEEkeywords}
biometrics; speaker recognition; feature extraction; feature matching and inference; active speaker detection; real-time computer vision; audio-visual fusion
\end{IEEEkeywords}

\section{Introduction}

Multimodal speaker attribution combining voice biometrics, face
re-identification, and active speaker detection (ASD) is generally
reported to outperform any single modality
\cite{chung2017lipsync,tao2021talknet,liao2023lightasd}. Almost all
such evaluations, however, report one aggregate figure per system,
computed over a benchmark whose composition (how many people are in
frame, how long the deployment has been running, how well individual
speakers are enrolled) is fixed by construction. A meeting-transcription
system deployed in practice sees none of that stability. Group size
varies from a one-on-one call to a seven-person meeting, and the same
system may run for months as its voice bank grows and its thresholds
get tuned. Whether the reported benefit of fusion holds uniformly
across these conditions, or is concentrated in some and nearly absent
in others, is a question that a single aggregate figure cannot answer.

We study this question using deployment logs from a real-time,
single-machine multimodal meeting tracker that attributes each speech
segment through an explicit \emph{priority cascade}, an ordered
sequence of decision rules in which the first rule whose evidence
requirement is met determines the final identity, rather than a single
learned fusion model. At its core, the underlying problem is a
multimodal \emph{biometric feature matching and inference} task, in
which voice and face embeddings are extracted independently by two
different feature-extraction pipelines (ECAPA-TDNN acoustic embeddings
and EdgeFace visual embeddings) and must be matched, cross-checked, and
reconciled online against a per-session identity registry. We use the
cascade's logged decisions, collected over 93 real, unscripted sessions
and 1{,}669 speech segments spanning 2.5 months, as an instrument to
measure how the necessity of each modality shifts across deployment
conditions rather than to report one more aggregate accuracy number.

Our contributions are the following.

\begin{enumerate}
\item Evidence that reliance on voice biometrics alone is highly
non-uniform across group sizes, ranging from 5\% to 40\% of segments,
with a structural explanation tied to which cascade rule takes
precedence at each group size, rather than a smooth trend.
\item A demonstration that an apparent decline in voice-only
resolution over a multi-month deployment is substantially confounded
by a shift in session composition, a methodological caution for
interpreting raw temporal trends in deployment logs without
controlling for this kind of shift.
\item A description of the rule-based audio-visual fusion cascade and
a lightweight self-diagnostic confidence-risk monitor that produced
the analyzed data, including its own real-time performance
characterization.
\item A discussion of what these findings imply for tuning and
evaluating fusion-based speaker attribution systems under variable
deployment conditions, and of the practical limitations of evaluating
such a system outside curated, ground-truth-labeled benchmarks.
\end{enumerate}

This work is part of a larger, ongoing multi-institution research
collaboration on multimodal meeting intelligence, of which the present
system is one of several parallel, complementary sub-projects; specific
institutional attributions are withheld here for double-blind review
and will be disclosed in the acknowledgments of the camera-ready
version.

\section{Related Work}

\textbf{Speaker diarization and verification.} pyannote.audio
\cite{bredin2020pyannote,bredin2023pyannote31} provides the diarization
backbone used in our pipeline, segmenting and clustering speech into
per-chunk speaker labels; diarization alone does not solve identity
persistence across non-adjacent processing windows. ECAPA-TDNN
\cite{desplanques2020ecapa}, trained on VoxCeleb, is the standard
architecture for the fixed-dimensional speaker embeddings we use both
for enrolled-speaker matching and for a persistent within-session
tracker for unenrolled speakers.

\textbf{Active speaker detection.} SyncNet-style lip-sync models
\cite{chung2017lipsync} and TalkNet \cite{tao2021talknet} established
the audio-visual correlation approach to ASD; Light-ASD
\cite{liao2023lightasd} is a lightweight CVPR 2023 model designed for
real-time inference, which we use as the primary ASD backend. We use
YOLOv8 \cite{jocher2023yolov8} for face detection, ByteTrack
\cite{zhang2022bytetrack} for short-term tracking, and EdgeFace
\cite{george2024edgeface} for lightweight face-recognition embeddings.

\textbf{Audio-visual fusion for speaker identity.} Most prior work on
audio-visual speaker identification targets offline, dataset-level
fusion, such as re-scoring diarization with ASD or jointly training
audio-visual embeddings, and reports a single accuracy figure over a
benchmark of fixed composition. Our setting differs in requiring an
\emph{online}, segment-by-segment decision under partial observability,
where a face, a voice match, or both may be missing or unreliable for
any given segment, which led us to a rule-based, explicitly ordered
fusion strategy rather than a single learned fusion score.

\textbf{Generalization across deployment conditions.} Benchmarks such
as AVA-ActiveSpeaker and VoxConverse fix group size and recording
conditions by construction, so a system's reported accuracy is, by
design, an average over whatever mixture of conditions the benchmark
happens to contain. We are not aware of prior work that stratifies a
deployed multimodal attribution system's behavior by group size or by
deployment time to test whether a single aggregate figure is
representative, which is the gap this paper addresses.

\section{System Overview}
\label{sec:system}

The system runs two cooperating threads sharing an in-memory state
object. A \emph{video thread} performs face detection (YOLOv8),
tracking (ByteTrack), face re-identification (EdgeFace), and per-frame
ASD updates. An \emph{audio thread} performs voice activity detection,
diarization-based session hints (pyannote.audio), transcription
(faster-whisper, \texttt{medium}, Portuguese), and speaker verification
(ECAPA-TDNN) on each detected speech segment, querying the video
thread's state for currently visible faces and ASD results. When
diarization indicates more than one concurrent speaker, the system
applies SepFormer \cite{subakan2021sepformer} speech separation before
re-verifying each separated stream.

Algorithm~\ref{alg:cascade} summarizes the fusion cascade that
attributes each segment. Each rule has a different evidence
requirement, and the cascade stops at the first rule whose precondition
is satisfied, preferring voice biometrics, which is invariant to camera
framing, whenever it is confident enough, and falling back to
progressively weaker, more circumstantial evidence otherwise. Two
cross-cutting passes run afterward. A \textbf{gender-consistency check}
rejects or reroutes visually derived decisions that contradict an
unambiguous audio-derived gender signal, skipped when voice-biometric
confidence is very high ($\geq 0.90$). A \textbf{turn-merging} pass
merges consecutive same-diarization-label segments only if their
voice-embedding similarity exceeds 0.40, since diarization occasionally
mislabels two different speakers under the same local cluster.

\begin{algorithm}[h]
\caption{Speaker attribution cascade (simplified)}
\label{alg:cascade}
\begin{algorithmic}[1]
\State $v \gets$ ECAPA verifier match (name, confidence) for segment
\State $g \gets$ audio-derived gender estimate
\State $a \gets$ ASD query over [segment start $-$ 0.4s, end $+$ 0.4s]
\If{$v$.conf $\geq 0.70$ \textbf{and} gender-consistent}
  \State \Return \textsc{Verifier-High} ($v$.name)
\ElsIf{$v$.conf $\geq 0.80$ \textbf{and} gender-consistent}
  \State \Return \textsc{Verifier-Mod} ($v$.name)
\ElsIf{$v$.conf $\geq 0.65$ \textbf{and} (ASD or gender corroborates)}
  \State \Return \textsc{Verifier-Weak} ($v$.name)
\ElsIf{$a$ is not \textsc{None}}
  \State \Return \textsc{ASD} (face bound to $a$)
\ElsIf{exactly 1 face visible, known identity, $\leq$1 expected speaker}
  \State \Return \textsc{Face-Only}
\ElsIf{exactly 1 face visible, raw\_conf $\geq 0.35$}
  \State \Return \textsc{Single-Face} (heuristic)
\ElsIf{ASD best-guess is dominant ($\geq$1.5--1.8$\times$ runner-up)}
  \State \Return \textsc{ASD-Guess}
\ElsIf{diarization session hint bound to a known face}
  \State \Return \textsc{Session-Face}
\ElsIf{diarization session hint available}
  \State \Return \textsc{Session-Hint} (anonymous)
\Else
  \State \Return \textsc{Session-Tracker} (persistent voice cluster)
\EndIf
\State \textit{(Gender cross-check / reroute applied as a final pass)}
\end{algorithmic}
\end{algorithm}

ASD itself uses a two-backend design, with Light-ASD
\cite{liao2023lightasd} as the primary backend and a model-free
pixel-difference fallback always available for windows Light-ASD has
not scored, for example immediately after model load or, in a subset
of our deployment sessions, when the neural backend failed to
initialize (\S\ref{sec:limitations}). The fallback rejects a common
source of false positives, occlusion, by requiring mouth-region motion
to exceed 1.8$\times$ the motion of the upper face before counting it
as speech-related.

\section{Experimental Setup}
\label{sec:setup}

We do not evaluate against a public diarization or ASD benchmark of
fixed composition, since doing so would reintroduce exactly the
single-condition averaging this paper argues against. Instead, we
instrument every real, unscripted use of the system over $\sim$2.5
months, from 16 March to 22 May 2026, logging for every processed
segment the final cascade decision type, the raw voice-verifier
confidence, whether ASD was active, end-to-end processing latency by
stage, a rule-based confidence-risk flag (\S\ref{sec:fprisk}), and the
session's configured number of expected speakers. The dataset comprises
93 recorded sessions and 1{,}669 processed speech segments, recorded in
Portuguese under variable lighting, off-camera speakers, overlapping
speech, partially enrolled participants, and consumer webcam and
microphone hardware, with between 1 and 7 expected speakers per
session. The system ran on a single machine (Intel Core i5-11400H,
8\,GB RAM, NVIDIA RTX 3060 Laptop GPU with 6\,GB VRAM) throughout the
deployment, so the latency figures reported below characterize this
one hardware configuration rather than a range of deployment targets;
a lower-end CPU-only machine would likely shift more segments above
the real-time threshold. Participants were members of the authors'
research group, informed that their meetings were being recorded to
develop and evaluate the system; recordings were not scripted or
staged for the study. We did not collect age, self-reported gender, or
other demographic attributes, so this dataset supports generalization
across the group-size and time conditions it was stratified by, not a
claim of demographic generalization.

\section{Results}

\subsection{Voice-only resolution is non-monotonic in group size}
\label{sec:bysize}

Averaged naively across the whole deployment, decisions based on pure
voice biometrics account for only 8.6\% of the 1{,}669 processed
segments, with the remaining 91.4\% relying on ASD, face
re-identification, gender disambiguation, or the diarization session
hint. That single number is already an argument against voice-only
attribution, but it is also, by construction, an average over sessions
that differed substantially in group size. Table~\ref{tab:byspeakers}
stratifies the same 1{,}669 segments by the
number of speakers configured for the session. Percentages below are
reported with 95\% Wilson confidence intervals, computed treating each
segment as an independent Bernoulli trial. The voice-only rate is
not a smooth function of group size, and each value has a structural
explanation rooted in the cascade itself (Algorithm~\ref{alg:cascade}),
not merely in how hard voice matching happens to be. Single-speaker
sessions show a low voice-only rate (9.9\%, 95\% CI [5.1, 18.3]) because
the cascade's
face-identity rules take precedence whenever at most one speaker is
expected, so voice matching is rarely the deciding rule even when it
would have agreed. Two-speaker sessions show the highest voice-only
rate observed (40.2\%, 95\% CI [31.7, 49.2]), plausibly because a
two-person voice bank is
easier to disambiguate and these sessions were disproportionately
smaller, better-enrolled meetings. Three-speaker sessions lean heavily
on the gender-consistency layer (23.3\%, 95\% CI [16.2, 32.3]) and, as
Table~\ref{tab:fprisk}
shows below, also carry the highest self-flagged risk of any group
size. Four-speaker sessions, which make up 82\% of the deployment,
drop to the lowest voice-only rate (5.0\%, 95\% CI [3.9, 6.3]) while ASD
carries the
largest share (36.8\%, 95\% CI [34.3, 39.4]). The wide intervals at the
smaller group sizes reflect their modest segment counts, but even the
loosest of them (three speakers, $n=103$) excludes both neighboring
group sizes' point estimates, so the pattern is not simply lost in
overlapping uncertainty. It shows that fusion is not
merely useful on average. It becomes more, not less, load-bearing as
group size grows toward the range most representative of real
meetings.

\begin{table}[h]
\centering
\caption{Decision source by number of expected speakers per session.}
\label{tab:byspeakers}
\footnotesize
\begin{tabular}{lrrrr}
\toprule
\textbf{Speakers} & \textbf{Segs.} & \textbf{Voice} & \textbf{ASD} & \textbf{Gender} \\
\midrule
1 &   81 &  9.9\% & 14.8\% &  1.2\% \\
2 &  117 & 40.2\% & 35.9\% &  0.9\% \\
3 &  103 & 19.4\% & 22.3\% & 23.3\% \\
4 & 1368 &  5.0\% & 36.8\% &  6.9\% \\
\bottomrule
\end{tabular}
\end{table}

Table~\ref{tab:fprisk} cross-checks this with the self-diagnostic risk
flag (\S\ref{sec:fprisk}). The risk rate is also non-monotonic and
peaks sharply at three speakers (68.9\%, 95\% CI [59.5, 77.1]), well
above both its
neighbors, matching the group size where the gender-safety layer had
to work hardest in Table~\ref{tab:byspeakers}. Three-speaker meetings
appear to be a structurally harder case for this cascade than either
smaller or larger groups, plausibly because they are large enough for
voice ambiguity to matter but small enough that gender alone often
fails to disambiguate among the candidates.

\begin{table}[h]
\centering
\caption{Self-flagged high-risk rate by number of expected speakers, with Wilson 95\% confidence intervals.}
\label{tab:fprisk}
\footnotesize
\begin{tabular}{lrrr}
\toprule
\textbf{Speakers} & \textbf{Segments} & \textbf{High-risk \%} & \textbf{95\% CI} \\
\midrule
1 &   81 & 21.0\% & [13.5, 31.1] \\
2 &  117 &  6.0\% & [\phantom{0}2.9, 11.8] \\
3 &  103 & 68.9\% & [59.5, 77.1] \\
4 & 1368 & 28.1\% & [25.8, 30.5] \\
\bottomrule
\end{tabular}
\end{table}

To test whether this non-monotonic pattern, and in particular the
three-speaker spike, is distinguishable from sampling noise, we ran a
chi-square test of homogeneity across the four group sizes (segment
level, treating decisions within a group as independent trials) for
both the voice-only rate ($\chi^2(3){=}187.4$, $p{=}2.2\times10^{-40}$)
and the high-risk rate ($\chi^2(3){=}113.6$, $p{=}1.8\times10^{-24}$),
and a two-sided Fisher exact test isolating the three-speaker cell
against the pooled remaining group sizes for the high-risk rate
(71/103 vs.\ 408/1{,}566, $p{=}5.0\times10^{-13}$) and the voice-only
rate (20/103 vs.\ 123/1{,}566, $p{=}3.3\times10^{-4}$). Both tests
reject the null hypothesis that group size is unrelated to either
rate, and the three-speaker cell alone is a significant outlier against
the pooled remainder, not merely against its immediate neighbors. We
flag one caveat these segment-level tests cannot resolve: segments
from the same session are not independent trials, and the
three-speaker cell's 103 segments come from only four sessions
(\S\ref{sec:limitations}), so its effective sample size for
generalizing to new three-speaker meetings is closer to four than to
103. The $p$-values above should be read as evidence that the spike is
not an artifact of within-session noise, not as a guarantee that its
exact magnitude would replicate over a larger set of three-speaker
sessions.

\subsection{Apparent temporal decline is a composition effect}
\label{sec:bytime}

A naive reading of the deployment log by calendar week shows what
looks like a steep decline in voice-only resolution. Voice-only
decisions accounted for 37--42\% of segments in the first two weeks,
falling to 0--2\% in the final three weeks recorded, alongside a
self-flagged risk rate that moved from an early peak near 52\% down to
14\% in the last week with sufficient volume to report. Read on its own, this trend invites a misleading
story, namely that voice matching degraded as the deployment
progressed.

Breaking down the same weeks by session composition rules this out.
The first two weeks were composed entirely of one- and two-speaker
sessions (100\% in each week), the group sizes Table~\ref{tab:byspeakers}
identifies as having the highest voice-only rates. From the third week
onward, the deployment shifted to sessions with three or four expected
speakers, reaching 90\% four-speaker segments by the third week and
staying between 82\% and 98\% for the remainder of the deployment. The
decline in voice-only resolution over time is, to a substantial degree,
the same non-monotonic group-size effect from
\S\ref{sec:bysize} appearing again under a different label, because the
deployment's session mix changed, not because voice biometrics got
worse at its job. The concurrent drop in self-flagged risk is
consistent with threshold and gender-check fixes made over the same
period (\S\ref{sec:limitations}), a genuine improvement running in
parallel with, and partly obscured by, the composition shift. We take
this as a methodological caution. A raw trend line over a real
deployment log can be confounded by what the deployment happened to
contain in each period, and stratifying by a plausible confounder,
here group size, before drawing conclusions from time alone is
necessary, not optional.

\subsection{Latency is not a confound}

A natural objection to the stratified findings above is that the
cascade might simply run slower, and therefore fall back to weaker
evidence more often, in harder conditions such as three-speaker
meetings. This is not the case: the real-time factor, total pipeline
latency divided by segment duration, was below 1 for 89.2\% of segments
overall (median 0.24), and the verifier, ASD query, and fusion logic
together add a median under 30\,ms regardless of condition, an order
of magnitude below transcription latency. The non-monotonic pattern in
\S\ref{sec:bysize} and \S\ref{sec:bytime} reflects which evidence was
available, not how fast the cascade could evaluate it.

\subsection{ASD backend accuracy against the published literature}
\label{sec:columbia}

We did not independently benchmark the ASD backend against a labeled
dataset (\S\ref{sec:limitations}). As a literature reference point, the
pretrained Light-ASD checkpoint we use is reported by its authors to
reach 94.06\% mAP on the AVA-ActiveSpeaker validation set and an
average F1 of 81.1\% on the Columbia ASD dataset
\cite{chakravarty2016columbia}, a small, five-panelist benchmark widely
used for out-of-domain ASD evaluation \cite{liao2023lightasd}, both
fixed-composition benchmarks of the kind this paper's deployment-scale
stratification is meant to complement.

\subsection{Self-diagnostic confidence-risk monitor}
\label{sec:fprisk}

Independently of the cascade, the system computes a rule-based risk
flag for every segment, based on factors such as a decision being of a
structurally weak type or an attributed name's biometric confidence
falling below a stricter 0.80 confirmation threshold even though it was
accepted by the cascade's looser threshold. This is not a ground-truth
accuracy measurement. It is a self-reported, deterministic function of
decision metadata intended to surface segments for human review without
requiring labeled data. Over the deployment, 479 of 1{,}669 segments
(28.7\%) were flagged high-risk overall, but as Table~\ref{tab:fprisk}
shows, this masks a range from 6.0\% to 68.9\% depending on group size,
reinforcing that a single risk-rate figure is as condition-dependent as
the decision-source figures it is meant to complement.

\section{Discussion and Limitations}
\label{sec:limitations}

\textbf{Design implications.} The stratified results suggest that
tuning a fusion cascade's thresholds against one representative
condition, or reporting one aggregate accuracy figure, risks missing
exactly the conditions where the system struggles most. Three-speaker
meetings, the sharpest outlier in both Table~\ref{tab:byspeakers} and
Table~\ref{tab:fprisk}, would be invisible in an evaluation that only
reported the 82\%-share four-speaker case or an unstratified average.
A practical implication is that cascade thresholds, or fallback
ordering, may need to be conditioned on expected group size rather than
fixed globally, an extension we have not implemented.

\textbf{No public benchmark evaluation, but a partial biased-sample
check.} We have not evaluated this
system on a labeled diarization or ASD benchmark, such as
AVA-ActiveSpeaker or VoxConverse, both of which fix composition by
construction in the way this paper argues against. We do, however,
have a small manually-corrected set from an earlier development
period, March to April 2026, before several threshold and
gender-check fixes: 24 sessions and 330 segments, selected during
development \emph{because} they were suspected to contain errors, not
sampled at random. Of the 124 segments where the cascade produced a
non-generic name, 37 (29.8\%, 95\% CI [22.5, 38.4]) matched the human
correction exactly; we do not report the corresponding 11.2\% figure
over all 330 segments (95\% CI [8.2, 15.1]) as representative of
current deployment accuracy, for the same error-selection reason and
because it predates later cascade fixes. What this biased sample can
support is a check on the risk monitor of \S\ref{sec:fprisk} itself:
among those 124 named predictions, correctness differed sharply by
risk flag, 81.8\% (9/11, 95\% CI [52.3, 94.9]) for unflagged versus
24.8\% (28/113, 95\% CI [17.7, 33.5]) for flagged, and a two-sided
Fisher exact test rejects independence between flag and correctness
($p=2.9\times10^{-4}$). This is not evidence the flag catches most
errors, since the sample was selected for containing them, but it is
evidence the flag is not noise: within this sample, a high-risk flag
came with roughly a three-fold drop in correctness. A properly
sampled, stratified ground-truth evaluation, covering all four group
sizes in Table~\ref{tab:byspeakers} and drawn without conditioning on
suspected errors, is the most important piece of future work here.

\textbf{Backend availability.} In a subset of deployment sessions the
Light-ASD backend failed to load due to a checkpoint key mismatch
between the released weights and the inference wrapper, triggering
fallback to the pixel-difference backend for that session's entire
duration. Our results are therefore a mixture of neural and
pixel-difference ASD, but which backend produced any individual ASD or
ASD-Guess decision (Algorithm~\ref{alg:cascade}) was not logged, only
the final decision type and confidence, so this is an instrumentation
gap rather than only a sample-size limitation. Logging backend
identity per decision is a low-cost fix that would let future work
quantify how much of the group-size and temporal patterns in
\S\ref{sec:bysize}--\S\ref{sec:bytime} are attributable to one backend
rather than the other.

\textbf{Session-level clustering behind the three-speaker cell.} The
103 segments in the three-speaker row of
Tables~\ref{tab:byspeakers}--\ref{tab:fprisk} are drawn from only four
sessions, so segments in that cell are not independent draws from the
population of three-speaker meetings in the way the segment-level
significance tests in \S\ref{sec:bysize} assume; a single unusual
session could dominate the cell. The Wilson intervals and $p$-values
reported there support that the spike is not sampling noise
\emph{within} the observed sessions, but a properly powered claim about
three-speaker meetings in general requires more independent sessions
at that group size, which is part of the same future-work gap noted
below for ground-truth evaluation.

\textbf{Diarization instability.} Per-chunk diarization labels are not
guaranteed to refer to the same physical speaker across two 8-second
processing windows. The session-hint and session-tracker fallbacks and
the voice-similarity merge gate (\S\ref{sec:system}) contain, but do
not eliminate, this instability.

\textbf{Threshold sensitivity and configuration drift.} The cascade's
decision boundaries rest on three manually tuned constants: the
verifier's own match threshold ($\tau=0.80$) and margin requirement
(0.07), which gate whether a name is returned at all, and a separate
confidence floor ($\tau_c=0.70$) that gates the
\textsc{Verifier-High} tier in Algorithm~\ref{alg:cascade}. Because a
name is only ever returned once it has already cleared $\tau$, and
$\tau_c<\tau$ under this deployment's standard configuration, the
\textsc{Verifier-Mod} tier is unreachable by construction: any segment
whose confidence clears $\tau_c=0.70$ has, by that point, already
cleared the stricter $\tau=0.80$ required to receive a name in the
first place, so \textsc{Verifier-High} always fires first and
\textsc{Verifier-Mod} never triggers. This is not a hypothetical
concern; it is what happened. Of the 92 sessions with a logged
configuration, 88 ran with $(\tau,\tau_c)=(0.80,0.70)$ and produced
zero \textsc{Verifier-Mod} decisions between them; all 16
\textsc{Verifier-Mod} segments in the deployment came from the
remaining 4 sessions, an early batch that instead ran $\tau=0.55$, a
looser value used before the threshold was tightened. Under that
older configuration $\tau_c=0.70$ sat \emph{above} $\tau$, so the tier
was reachable for confidences in $[0.55,0.70)$ and disappeared as soon
as $\tau$ was raised past $\tau_c$ later in the deployment. An
ordinary manual retuning of one threshold silently orphaned an entire
cascade tier for the remainder of the deployment, with no error and no
conspicuous accuracy drop, since \textsc{Verifier-High} simply
absorbed the segments \textsc{Verifier-Mod} would otherwise have
handled. This is a concrete instance of the general risk the review
raised: fixed, manually tuned thresholds in a multi-tier cascade can
interact so that a rule is dead code under one configuration and live
under another, and nothing in the cascade's own output flags which
regime is active without checking the boundary condition directly, as
we did here. A sensitivity check that varies each threshold
independently and reports which tiers remain reachable, run
automatically whenever a threshold changes, would catch this class of
issue before deployment rather than after the fact. Two other drift
sources remain unevaluated: sensor drift (webcam or microphone
characteristics or placement changing across sessions, not logged in
a way we can check) and concept drift proper (enrolled speakers'
voices or appearance changing gradually over the 2.5-month deployment,
against which our identity registry has no mechanism to adapt); both
are plausible additional contributors to the temporal patterns in
\S\ref{sec:bytime} that our composition-shift explanation does not
rule out.

\textbf{Rule-based vs. learned fusion.} The cascade's thresholds were
tuned manually against observed failures rather than learned from
labeled data, which keeps decisions auditable but does not adapt
automatically to a shifting condition mix like the one documented in
\S\ref{sec:bytime}. A learned, condition-aware re-ranking of cascade
outputs is a natural extension once labeled data spanning multiple
group sizes is available.

\textbf{Generalization beyond meetings.} The cascade's tier ordering,
voice biometrics first, then active-speaker cues, then face and
gender agreement, then session context, was tuned for small, seated
meetings recorded with a mostly frontal consumer webcam. We expect the
\emph{stratification methodology}, not any specific threshold value in
Algorithm~\ref{alg:cascade}, to transfer to other domains, but the
cascade itself would likely shift tiers differently in each. In
classrooms or lecture halls, wider and less speaker-frontal camera
framing and a mostly unenrolled voice roster would push more decisions
toward the ASD and session-context tiers, similar to this deployment's
four-speaker case. In customer-service or reception settings, a small,
well-enrolled staff roster and frequent one-speaker turns resemble our
best-performing one- and two-speaker cases and would likely sustain a
higher voice-only rate. Surveillance-style deployments, with
uncooperative or unenrolled subjects and oblique camera angles, would
strain the face-identity and gender-consistency tiers the most, since
both assume a cooperative, enrolled subject that such a deployment
cannot guarantee; a domain like that would need to lean more heavily
on voice and active-speaker evidence, or accept a much higher
session-context (anonymous) rate than reported here.

\section{Conclusion}

Aggregate deployment statistics can conceal as much as they reveal. In
93 real meeting sessions and 1{,}669 speech segments, voice biometrics
alone resolved 8.6\% of segments overall, but that single number hid a
five-fold range, from 5\% to 40\%, depending on expected group size,
and an apparent time trend that was largely a composition effect in
disguise. The priority cascade that generated this data sustains
real-time throughput and flags its own uncertainty without labeled
data, but the more general lesson is methodological. Reporting or
tuning a multimodal attribution system against one aggregate figure
risks hiding exactly the conditions, here three-speaker meetings,
where it is least reliable. Stratifying deployment logs by plausible
confounders such as group size, before trusting any trend observed
over time, is a cheap and necessary check that a fixed-composition
benchmark cannot substitute for. The most important next step is a
properly sampled, ground-truth-labeled evaluation stratified the same
way, to replace the deployment-distribution statistics reported here
with calibrated accuracy figures per condition.

\bibliographystyle{IEEEtran}
\bibliography{references}

\end{document}
